% Präambel
\documentclass[12pt,			% Schriftgröße
a4paper,						% Papierformat
twoside, 						% einseitiges (oneside) oder zweiseitiges (twoside) Dokument
listof=totoc, 					% Tabellen- und Abbildungsverzeichnis ins Inhaltsverzeichnis
bibliography=totoc,				% Literaturverzeichnis ins Inhaltsverzeichnis aufnehmen
titlepage, 						% Titlepage-Umgebung statt \maketitle
headsepline, 					% horizontale Linie unter Kolumnentitel
%abstracton,					% Überschrift beim Abstract einschalten, Abstract muss dazu in {abstract}-Umgebung stehen
DIV18,							% auskommentieren, um den Seitenspiegel zu vergrößern
BCOR6mm,						% Bindekorrektur, die den Seitenspiegel um 6mm nach rechts verschiebt,
cleardoublepage=empty,			% Stil einer leeren eingefügten Seite bei Kapitelwechsel
parskip,						% Absatzabstand bei Absatzwechsel einfügen
ngerman							% Übersetzung der festcodierten Überschriften
]{scrbook}			
\usepackage{ucs} 				% Dokument in utf8-Codierung schreiben und speichern
\usepackage{babel} 				% deutsche Trennungsregeln
\usepackage[utf8x]{inputenc} 	% ermöglicht die direkte Eingabe von Umlauten
\usepackage[T1]{fontenc} 		% Ausgabe aller zeichen in einer T1-Codierung (wichtig für die Ausgabe von Umlauten!)
\setlength{\parindent}{0ex} 	% bei neuem Abschnitt nicht einrücken
\linespread{1.2}\selectfont     % Zeilenabstand erhöhen - größere Werte als 1.2 nicht verwenden!!
\usepackage{scrpage2}			% SCR Headings verwenden
\setheadsepline{0.4pt}			% Kopfzeile Linien oben
\setfootsepline{0.4pt}			% Kopfzeile Linien unten
\pagestyle{scrheadings}			% SCR Headings einschalten
\usepackage{graphicx}  			% Einbinden von Grafiken erlauben
\usepackage[format=hang,		% Formatierungen von Unter- / Überschriften
font=normal,
labelfont=bf,
justification=RaggedRight,
singlelinecheck=true,
aboveskip=1mm
]{caption}

\usepackage{enumitem}			% Erlaubt Änderung der Nummerierung in der Umgebung enumerate

\usepackage{amsmath}			% Ergänzungen für Formeln
\usepackage{textcomp} 			% zum Einsatz von Eurozeichen u. a. Symbolen
\usepackage[%
%per=slash,
decimalsymbol=comma,
%thickspace,					
%thickqspace
]{siunitx}						% Einfache Eingabe von SI-Einheiten
\sisetup{number-unit-product = \;, inter-unit-product = \:}
\newcommand*\diff{\mathop{}\!\mathrm{d}}	% Differentialzeichen
\newcommand*\Diff[1]{\mathop{}\!\mathrm{d^#1}} % Differentialzeichen höherer Ableitung
\newcommand*\jj{\mathop{}\!\mathrm{j}}	% Komplexe Zahl j

\usepackage[					% Einstellunge Paket hyperref
hyperfootnotes=false			% im pfd-Output Fußnoten nicht verlinken
]{hyperref}

\usepackage{makeidx}			% Paket zur Erstellung eines Index
\usepackage[printonlyused]{acronym} 	% zur Erstellung des Abkürzungsberzeichnisses

\usepackage[					% Einstellungen für Fußnoten
bottom,							% Ausrichtung unten
multiple,						% Trennung durch Seperator bei mehreren Fußnoten
hang,
marginal
]{footmisc}

\usepackage{calc}				% Paket zum Berechnen von Längen z.B. 0.8\linewidth

\usepackage{xcolor} 			% einfache Verwendung von Farben in nahezu allen Farbmodellen

\usepackage{listings}			% Darstellung von Quellcode mit den Umgebungen {lstlisting}, \lstinline und \lstinputlisting
\lstset{literate=				% Damit können Umlaute innerhalb Listings geschrieben werden
	{Ö}{{\"O}}1
	{Ä}{{\"A}}1
	{Ü}{{\"U}}1
	{ß}{{\ss}}1
	{ü}{{\"u}}1
	{ä}{{\"a}}1
	{ö}{{\"o}}1
}

%****************************************************************** Zusätzliche Packages *******************
%*****************************************************

\usepackage{color}
\usepackage{colortbl}
\usepackage{pdfpages}

%*****************************************************************************************************************************************************************

\definecolor{mygreen}{rgb}{0,0.6,0}
\definecolor{mygray}{rgb}{0.5,0.5,0.5}
\definecolor{mymauve}{rgb}{0.58,0,0.82}
\lstset{ %
	backgroundcolor=\color{white},   % choose the background color; you must add \usepackage{color} or \usepackage{xcolor}; should come as last argument
	basicstyle=\footnotesize,        % the size of the fonts that are used for the code
	breakatwhitespace=false,         % sets if automatic breaks should only happen at whitespace
	breaklines=true,                 % sets automatic line breaking
	captionpos=t,                    % sets the caption-position to (b) bottom or (t) top
	commentstyle=\color{mygreen},    % comment style
	deletekeywords={...},            % if you want to delete keywords from the given language
	escapeinside={\%*}{*)},          % if you want to add LaTeX within your code
	escapeinside={(*@}{@*)},
	extendedchars=true,              % lets you use non-ASCII characters; for 8-bits encodings only, does not work with UTF-8
	frame=none,	                   	% "single" adds a frame around the code; "none"
	keepspaces=true,                 % keeps spaces in text, useful for keeping indentation of code (possibly needs columns=flexible)
	keywordstyle=\color{blue},       % keyword style
	language=[LaTeX]TeX,             % the language of the code
	morekeywords={*,nomenclature},   % if you want to add more keywords to the set
	numbers=left,                    % where to put the line-numbers; possible values are (none, left, right)
	numbersep=5pt,                   % how far the line-numbers are from the code
	numberstyle=\tiny\color{mygray}, % the style that is used for the line-numbers
	rulecolor=\color{black},         % if not set, the frame-color may be changed on line-breaks within not-black text (e.g. comments (green here))
	showspaces=false,                % show spaces everywhere adding particular underscores; it overrides 'showstringspaces'
	showstringspaces=false,          % underline spaces within strings only
	showtabs=false,                  % show tabs within strings adding particular underscores
	stepnumber=1,                    % the step between two line-numbers. If it's 1, each line will be numbered
	stringstyle=\color{mymauve},     % string literal style
	tabsize=2,	                   % sets default tabsize to 2 spaces
	title=\lstname                   % show the filename of files included with \lstinputlisting; also try caption instead of title
}

\makeindex						% Indexverzeichnis erstellen
				% Abkürzungsverzeichnis erstellen

% -----------------------------------------------------------------------------------------------------------------
% Zum Aktualisieren des Abkürzungsverzeichnisses (Nomenklatur) bitte auf der Kommandozeile folgenden Befehl aufrufen :
% makeindex <Dateiname>.nlo -s nomencl.ist -o <Dateiname>.nls
% Oder besser: Kann in TexStudio unter Tools-Benutzer als Shortlink angelegt werden
% Konfiguration unter: Optionen-Erzeugen-Benutzerbefehle: makeindex -s nomencl.ist -t %.nlg -o %.nls %.nlo
% -----------------------------------------------------------------------------------------------------------------

% Hier die persönlichen Daten eingeben:

\newcommand{\titel}{Eintwicklung eines CO2-Messers}
\newcommand{\untertitel}{}
\newcommand{\arbeit}{Mikrocomputertechnik - Bericht}
\newcommand{\studiengang}{Elektrotechnik}
\newcommand{\studienrichtung}{Fahrzeugelektronik}
\newcommand{\autor}{Alexander Herrmann, Johannes Ruffer, Serkant Soylu}
\newcommand{\matrikelnr}{9859538 x 1011921 x 9964027}
\newcommand{\kurs}{TFE18-2}
\newcommand{\abgabe}{19.04.2020}
\newcommand{\zeitraum}{01.10.2019 - 19.04.2020}
\newcommand{\betreuerdhbw}{Hans Jürgen Herpel}
\newcommand{\jahr}{2020}			% für Angabe im Copyright-Vermerk der Titelseite

% Folgende Zeilen definieren Abkürzungen, um Befehle schneller eingeben zu können
\newcommand{\ua}{\mbox{u.\,a.\ }}
\newcommand{\zB}{\mbox{z.\,B.\ }}
\newcommand{\bs}{$\backslash$}

% Folgende Zeilen weden benötigt, um Tikz und PGF-Plot-Grafiken einzubinden
\usepackage{pgfplots}
\usepackage{pgfplotstable}
\pgfplotsset{compat=newest,width=0.6\linewidth}
\usepgfplotslibrary{smithchart}
\usepackage{tikz}						% Tikz sollte nach Listings Pakete geladen werden.
\usetikzlibrary{arrows}
\usetikzlibrary{positioning,shadings}


%\usepackage{pgfkeys,pgfmath,pgfcore}% or pgf or tikz or pgfplots
%
%\pgfkeys{
%	/pgf/number format/textnumber/.style={
%		fixed,
%		use comma,
%		fixed zerofill,
%		precision=4,
%		1000 sep={.},
%	},
%}


\hyphenation{Schrift-ar-ten}

%\renewcommand\thesection{~\arabic{section}}



% -------------------------------------------------------------------------------------------
%                     Beginn des Dokumenteninhalts
% -------------------------------------------------------------------------------------------
\begin{document}
\setcounter{secnumdepth}{3}				% Nummerierungstiefe fürs Inhaltsverzeichnis
\setcounter{tocdepth}{3}
\sffamily								% für die Titelei serifenlose Schrift verwenden

% ------------------------------ Titelei -----------------------------------------------------

\thispagestyle{plain}
\begin{titlepage}
\enlargethispage{4.0cm}
\sffamily 								% Serifenlose Grundschrift für die Titelseite einstellen

\parbox{0.5\linewidth}{
\begin{flushleft}
%\includegraphics[width=0.4\linewidth]{images/Daimler-Logo}\\[5ex]
\end{flushleft}
}
\parbox{0.5\linewidth}{
\begin{flushright}
	\includegraphics[width=0.4\linewidth]{images/DHBW_d_R_FN_46mm_4c}\\[5ex]
\end{flushright}
}
				

\begin{center}

\LARGE{\textsc{\textbf{\titel}}}\\[1.5ex]
%\Large{\textbf{\untertitel}}\\[5ex]
\Large{\textbf{\arbeit}}\\[2ex]
%\normalsize{für die Prüfung zum\\[1ex] Bachelor of Engineering}\\[3ex]
\Large{Studiengang \studiengang}\\[2ex]
\large{Studienrichtung \studienrichtung}\\[1ex]
\normalsize{Duale Hochschule Baden-Württemberg Ravensburg, Campus Friedrichshafen}\\[5ex]
von\\[1ex] 

\begin{tabular}{ccc}
	& & \\
	Alexander Herrmann & Johannes Ruffer & Serkant Soylu \\
\end{tabular}


\end{center}

\begin{flushleft}

\begin{tabular}{ll}
Abgabedatum:					& \quad \abgabe \\
Bearbeitungszeitraum:		   		& \quad \zeitraum \\ 
Matrikelnummer: 			& \quad \matrikelnr \\
Kurs: 							& \quad \kurs \\
Gutachter der Dualen Hochschule: & \quad \betreuerdhbw \\ [5ex]

\end{tabular} 



\small
%Copyrightvermerk:\\

%Dieses Werk einschließlich seiner Teile ist \textbf{urheberrechtlich geschützt}. Jede Verwertung außerhalb der engen Grenzen des Urheberrechtgesetzes ist ohne Zustimmung des Autors unzulässig und strafbar. Das gilt insbesondere für Vervielfältigungen, Übersetzungen, Mikroverfilmungen sowie die Einspeicherung und Verarbeitung in elektronischen Systemen.
\end{flushleft}
\begin{flushright}
%\copyright{} \jahr
\end{flushright}
\end{titlepage}

\cleardoublepage 				% erzeugt die Titelseite
\pagenumbering{roman}					% kleine, römische Seitenzahlen für Titelei
\chapter*{Eidesstattliche Erklärung} %*-Variante sorgt dafür, das Abstract nicht im Inhaltsverzeichnis auftaucht

%Mit der Abgabe der Aufgabenlösung in Moodle erklären Sie, dass Sie den Programmentwurf selbständig verfasst und keine anderen als die angegebenen Quellen und Hilfsmittel benutzt haben.

Gemäß Ziffer 1.1.13 der Anlage 1 zu §§ 3, 4 und 5  der Studien- und Prüfungsordnung für die Bachelorstudiengänge im Studienbereich Technik der Dualen Hochschule Baden-Würt­tem­berg vom 29.09.2015.

Wir versichern hiermit, dass wir unsere Projektarbeit mit dem Thema: 

\begin{quote}
	\textit{\titel} \textit{ \untertitel }
\end{quote}

selbstständig verfasst und keine anderen als die angegebenen Quellen und Hilfsmittel benutzt haben. Wir versichern zudem, dass die eingereichte elektronische Fassung mit der gedruckten Fassung übereinstimmt.

Friedrichshafen, den \today \\[2ex]

\begin{center}

\begin{tabular}{ccc}
	\rule[-0.2cm]{0.3\linewidth}{0.5pt} & \rule[-0.2cm]{0.3\linewidth}{0.5pt} & \rule[-0.2cm]{0.3\linewidth}{0.5pt}\\
	Alexander Herrmann & Johannes Ruffer & Serkant Soylu\\
\end{tabular}

\rule{0.75\linewidth}{0.1pt}\\
\textsc{Autoren} \\[10ex]
\end{center}


%Sperrvermerk bei Bedarf dekommentieren
%\hrule 
%\vspace*{1.0cm}
%\noindent \textbf{\Large{Sperrvermerk}}\\
%\normalsize
%Die Ergebnisse der Arbeit stehen ausschließlich dem auf dem Deckblatt aufgeführten Ausbildungsbetrieb zur Verfügung.

\normalsize

\rmfamily



\cleardoublepage 				% Einbinden der eidestattlichen Erklärung

\addchap*{Kurzfassung}
% Kurzfassung - Autor:

\label{kurzfassung}

Der Bericht wurde von drei Studierenden an der \ac{DHBW}-Ravensburg Campus Friedrichshafen im Rahmen der Mikrocomputertechnik Vorlesung eigenständig von der Projektplanung, über die Durchführung, bis hin zum Projektabschluss mit Dokumentation durchgeführt. \\
Ziel der Arbeit ist es, durch die Praxiserfahrung mit Mikrocomputern, Fähigkeiten und Wissen in diesem Bereich zu erwerben. Durch das Vergleichen von anfänglichen Kosten- und Komplexitätsschätzungen mit den späteren Ergebnissen in der Umsetzung können die Studierenden ein Fazit ziehen, inwiefern diese übereinstimmen. So werden auch Fähigkeiten im Bereich des Projektmanagements weiterentwickelt. \\
Dieses Projekt befasst sich mit der Simulation eines Fensterhebers. Dabei wird das Öffnen und Schließen der Fenster mithilfe einer \ac{LED} simuliert, welche nach der Überschreitung eines bestimmten Grenzwertes aufleuchtet. Sobald nach ausreichendem Lüften die Luftgüte unter einen Schwellwert gerät, soll die \ac{LED} wieder aus gehen. \\
Damit dies möglich ist, wurde die Entwicklung von Software-Code, sowie ein Schaltungslayout und der 3D-Druck des Gehäuses selbständig vorgenommen. \\
Der folgende Bericht dokumentiert die Motivation, Vorgehensweise und Umsetzung von der Projektplanung, über die Implementierung von Hardware und Software, bis hin zu durchgeführten Tests und dem Projektabschluss.

\tableofcontents						% Erzeugen des Inhalsverzeichnisses
\cleardoublepage

% --------------------------------------------------------------------------------------------
%                    Inhalt der Bachelorarbeit
%---------------------------------------------------------------------------------------------
\pagenumbering{arabic}					% arabische Seitenzahlen für den Hauptteil

\rmfamily

\chapter{Einleitung}
%\setcounter{figure}{0}
\setcounter{chapter}{1}
% Autor:

\label{Einleitung}

Mithilfe eines Arduinos wurde in diesem Projekt ein CO2-gesteuerter Fensterheber simuliert, welcher dazu dienen soll die Räumlichkeiten bei schlechter Luftqualität automatisch zu lüften. Auch das automatische schließen des Fensters nach Wiederherstellung von guter Luftgüte wird simuliert. \\
Das Öffnen und Schließen der Fenster wird anhand von einer LED simuliert, welche nach der Überschreitung eines bestimmten Grenzwertes aufleuchtet. Sobald nach ausreichendem Lüften die Luftgüte unter einen Schwellwert gerät, soll die LED wieder aus gehen. \\
Damit dies möglich ist, wurde die Entwicklung von Software-Code, sowie ein Schaltungslayout und der 3D-Druck des Gehäuses selbständig vorgenommen. \\
Der folgende Bericht dokumentiert die Vorgehensweisen und Umsetzung von der Projektplanung, über die Implementierung von Hardware und Software, bis hin zu durchgeführten Tests und den Projektabschluss.

\chapter{Erläuterung der Aufgabenstellung}
\setcounter{chapter}{2}
% Autor:

\label{Aufgabenstellung}




\chapter{Ausblick}
\label{Ausblick}

In dem folgenden Abschnitt ist der Ausblick, ebenso wie das Fazit, aufgrund der parallelen Entwicklung während des Projektes, in zwei Teile gegliedert. Der Ausblick soll zeigen, welche Möglichkeiten es gibt, dieses Projekt weiterzuentwickeln und zu optimieren.

\chapter {Fazit}
%Fazit - Autor:
\label{Fazit}



\addchap{Autorenverzeichnis}
\renewcommand\thesection{\arabic{section}}
Alexander Herrmann\\
Johannes Ruffer\\
Serkant Soylu\\

\addchap{Verzeichnis verwendeter Abkürzungen und Formelzeichen}

\begin{acronym}%[CSMA/CD + AMP]
	
	\acro{Bsp}{Beispiel}
	
\end{acronym}

% ---- Literaturverzeichnis ----------

\bibliography{Literatur/literatur}
										% Einbindung mehrerer Verzeichnisse in einem \bibliography Befehl mit Kommata trennen - keine Leerzeichen nach den Kommata!
\bibliographystyle{alphadin}           	%plain: alphabetisch, unsrt: nach Zitat, alphadin: NameJahr


% -----Ausgabe aller Verzeichnisse ---
\setlength{\parskip}{0.5\baselineskip}
\renewcommand{\indexname}{Sachwortverzeichnis}
\printindex								% Erzeugen des Indexverzeichnises
\addcontentsline{toc}{chapter}{\indexname}
% Version 0.1 vom 31.10.14 - T. Kibler

% Änderungswünsche bitte an T. Kibler melden

% alle Abkürzungen, die in der Arbeit verwendet werden. Nachfolgend sind allgemeine Abkürzungen aufgelistet. Die für die Arbeit spezifischen Abkürzungen sollten innerhalb des Haupttextes aufgeführt werden. Die Alphabetische Sortierung übernimmt Latex.

				% Datei mit allgemeinen Abkürzungen laden
%\renewcommand{\nomname}{Verzeichnis verwendeter %Formelzeichen und Abkürzungen}
%\setlength{\nomlabelwidth}{.20\hsize}
%\renewcommand{\nomlabel}[1]{#1 \dotfill}
%\setlength{\nomitemsep}{-\parsep}
%\printnomenclature						% Erzeugen des Abkürzungsverzeichnises, siehe auch Inhalt der Datei pages/abkuerzungen.tex
%\cleardoublepage
%\renewcommand{\glossaryname}{Glossar}
%\printglossaries
%\cleardoublepage
\listoffigures 							% Erzeugen des Abbildungsverzeichnisses 
\cleardoublepage
\listoftables 							% Erzeugen des Tabellenverzeichnisses
\cleardoublepage

% -----Anhang ------------------------
\appendix
\chapter{Anhang}

\section{Weitere Abbildungen}
%\clearpage
%\pagenumbering{Roman}					% große, römische Seitenzahlen für Anhang, falls gewünscht

\end{document}